% Generated by the blueprint tool (blueprint-pdf/1) from a blueprint snapshot; edit the
% Markdown sources, not this file.  Compile with `latexmk -pdf` or by running
% `pdflatex` twice (the second run settles the contents and the page references).
% Images, if any, are read from assets/ next to this file.
\documentclass[11pt]{article}
\usepackage[T1]{fontenc}
\usepackage[utf8]{inputenc}
\usepackage{lmodern}
\usepackage{textcomp}
\usepackage[a4paper,margin=2.5cm,headheight=14pt]{geometry}
\usepackage{amsmath,amssymb,amsthm}
\usepackage{amscd}  % KaTeX's CD diagrams; part of amsmath, so always there
\usepackage{mathtools}
\usepackage{bm}
\usepackage{graphicx}
\usepackage{xcolor}
\usepackage{array,booktabs,tabularx}
\usepackage{fancyvrb}
\usepackage[normalem]{ulem}
\usepackage{microtype}
% Optional, and only with their fonts: a small TeX installation has the
% packages without the fonts, which fails at the first formula.
\IfFileExists{mathrsfs.sty}{\IfFileExists{rsfs10.tfm}{\usepackage{mathrsfs}}{}}{}
\IfFileExists{stmaryrd.sty}{\IfFileExists{stmary10.tfm}{\usepackage{stmaryrd}}{}}{}
\IfFileExists{cancel.sty}{\usepackage{cancel}}{}
\usepackage[hyphens]{url}
\usepackage[pdfusetitle,bookmarksnumbered=false]{hyperref}
\usepackage{bookmark}
\definecolor{bplink}{HTML}{1F4E9C}
\definecolor{bpmuted}{HTML}{5F676E}
\hypersetup{colorlinks=true,linkcolor=bplink,urlcolor=bplink,citecolor=bplink}
\setlength{\parindent}{0pt}
\setlength{\parskip}{0.55em plus 0.2em minus 0.1em}
% Lists as tight as the page's.
\makeatletter
\def\@listi{\leftmargin\leftmargini\parsep 0pt\topsep 0.3em plus 0.1em\itemsep 0.1em plus 0.1em}
\let\@listI\@listi
\def\@listii{\leftmargin\leftmarginii\labelwidth\leftmarginii\advance\labelwidth-\labelsep\parsep 0pt\topsep 0.1em\itemsep 0.1em}
% An ordered list that does not start at 1.
\newcommand{\bpstartat}[1]{\setcounter{\@enumctr}{\numexpr#1-1\relax}}
\makeatother
\setcounter{tocdepth}{3}
\emergencystretch=3em
\allowdisplaybreaks

% A numbered statement, set as a paper sets one: "Theorem 2.1 (Title)."
% #1 is the lead ("Theorem~2.1"), #2 what follows it in roman.
\newenvironment{bpstatement}[2]{\par\addvspace{0.6\baselineskip}\noindent\textbf{#1}#2\textbf{.}\enspace\ignorespaces}
  {\par\addvspace{0.4\baselineskip}}
\newenvironment{bpstatementit}[2]{\begin{bpstatement}{#1}{#2}\itshape\ignorespaces}{\end{bpstatement}}
% A step: a relation with prose of its own, read under the object it starts from.
\newenvironment{bpstep}[1]{\par\addvspace{0.3\baselineskip}\begin{list}{}{\leftmargin=1.6em\rightmargin=0pt\topsep=0pt}\item\relax
  {\small\color{bpmuted}#1}\par\nopagebreak}{\end{list}\par}
% Notes about where an object is written out.
\newenvironment{bpnote}{\par\small\color{bpmuted}\itshape}{\par}
\newcommand{\bpomitted}[1]{\par{\small\color{bpmuted}\itshape #1\par}}
\newcommand{\bpbodyheading}[1]{\par\addvspace{0.4\baselineskip}\noindent\textbf{#1}\par\nopagebreak}
\newcommand{\bprule}{\par\medskip\noindent\hfil\rule{0.3\linewidth}{0.4pt}\hfil\par\medskip}
\newcommand{\bpboundary}[1]{\par\nopagebreak{\small\color{bpmuted}#1\par}\nopagebreak}
% A reference to something this export leaves out, and one that resolves nowhere.
\newcommand{\bpout}[1]{#1\textsuperscript{\hyperlink{bp-outside}{\ensuremath{\ast}}}}
\newcommand{\bpbroken}[1]{\textcolor{red!70!black}{[#1?]}}
\newcommand{\bpslug}[1]{\texttt{#1}}
% A bibliography entry: label, anchor, entry; the label hangs in the margin.
\newcommand{\bpbibitem}[3]{\par\setbox0\hbox{[#1]\enspace}\ifdim\wd0<5em\setbox0\hbox to 5em{[#1]\hfil}\fi
  \hangindent=\wd0 \hangafter=1 \noindent\box0\phantomsection\label{#2}#3\par\smallskip}
% A character the export knows no way to set.
\newcommand{\bpglyph}[1]{\fbox{\scriptsize #1}}
% Images, at their own size unless they are wider than the line.
\newsavebox{\bpimagebox}
\newcommand{\bpimage}[1]{\sbox{\bpimagebox}{\includegraphics{#1}}%
  \ifdim\wd\bpimagebox>\linewidth\resizebox{\linewidth}{!}{\usebox{\bpimagebox}}\else\usebox{\bpimagebox}\fi}
\newcommand{\bpmissingimage}[1]{\fbox{\parbox{0.8\linewidth}{\centering\small\color{bpmuted}Image not included: #1}}}
\DefineVerbatimEnvironment{bpverbatim}{Verbatim}{fontsize=\small,xleftmargin=1em}
\renewcommand{\qedsymbol}{\ensuremath{\blacksquare}}
\hypersetup{pdftitle={Weil Conjectures — F Algebraic groups in characteristic zero},pdfsubject={Appendix F Algebraic groups in characteristic zero; all levels},pdfcreator={blueprint}}
% KaTeX commands the website knows
\providecommand{\R}{\mathbb{R}}\providecommand{\N}{\mathbb{N}}\providecommand{\Z}{\mathbb{Z}}
\providecommand{\Q}{\mathbb{Q}}\providecommand{\C}{\mathbb{C}}
\providecommand{\Reals}{\mathbb{R}}\providecommand{\reals}{\mathbb{R}}\providecommand{\Complex}{\mathbb{C}}
\providecommand{\cnums}{\mathbb{C}}\providecommand{\natnums}{\mathbb{N}}\providecommand{\Bbb}[1]{\mathbb{#1}}
\providecommand{\bold}[1]{\mathbf{#1}}\providecommand{\frak}[1]{\mathfrak{#1}}\providecommand{\mathbbm}[1]{\mathbb{#1}}
\providecommand{\lang}{\langle}\providecommand{\rang}{\rangle}\providecommand{\bra}[1]{\langle #1\rvert}
\providecommand{\ket}[1]{\lvert #1\rangle}\providecommand{\braket}[1]{\langle #1\rangle}
\providecommand{\set}[1]{\{#1\}}\providecommand{\Set}[1]{\left\{#1\right\}}
\providecommand{\infin}{\infty}\providecommand{\isin}{\in}\providecommand{\plusmn}{\pm}\providecommand{\sub}{\subset}
\providecommand{\sube}{\subseteq}\providecommand{\supe}{\supseteq}\providecommand{\exist}{\exists}
\providecommand{\bull}{\bullet}\providecommand{\larr}{\leftarrow}\providecommand{\rarr}{\rightarrow}
\providecommand{\lrarr}{\leftrightarrow}\providecommand{\Larr}{\Leftarrow}\providecommand{\Rarr}{\Rightarrow}
\providecommand{\Lrarr}{\Leftrightarrow}\providecommand{\uarr}{\uparrow}\providecommand{\darr}{\downarrow}
\providecommand{\harr}{\leftrightarrow}\providecommand{\lArr}{\Leftarrow}\providecommand{\rArr}{\Rightarrow}
\providecommand{\hArr}{\Leftrightarrow}\providecommand{\alef}{\aleph}\providecommand{\alefsym}{\aleph}
\providecommand{\real}{\Re}\providecommand{\image}{\Im}\providecommand{\weierp}{\wp}
\providecommand{\Alpha}{\mathrm{A}}\providecommand{\Beta}{\mathrm{B}}\providecommand{\Epsilon}{\mathrm{E}}
\providecommand{\Zeta}{\mathrm{Z}}\providecommand{\Eta}{\mathrm{H}}\providecommand{\Iota}{\mathrm{I}}
\providecommand{\Kappa}{\mathrm{K}}\providecommand{\Mu}{\mathrm{M}}\providecommand{\Nu}{\mathrm{N}}
\providecommand{\Omicron}{\mathrm{O}}\providecommand{\Rho}{\mathrm{P}}\providecommand{\Tau}{\mathrm{T}}
\providecommand{\Chi}{\mathrm{X}}\providecommand{\omicron}{o}\providecommand{\gt}{>}\providecommand{\lt}{<}
\providecommand{\And}{\&}\providecommand{\hdashline}{\hline}\providecommand{\operatornamewithlimits}[1]{\operatorname*{#1}}
\providecommand{\llbracket}{[\![}\providecommand{\rrbracket}{]\!]}\providecommand{\cancel}[1]{#1}
\providecommand{\mathscr}{\mathcal}
% the characters beyond ASCII this document uses
\DeclareUnicodeCharacter{2162}{\textup{III}}
% the project's macros ([katex.macros] in blueprint.toml)
\providecommand{\A}{}\renewcommand{\A}{\mathbf A}
\providecommand{\AEn}{}\renewcommand{\AEn}{A_{E,n}}
\providecommand{\Aaf}{}\renewcommand{\Aaf}{\mathbf A}
\providecommand{\Ab}{}\renewcommand{\Ab}{\mathsf{Ab}}
\providecommand{\AffSch}{}\renewcommand{\AffSch}{\operatorname{AffSch}}
\providecommand{\Aone}{}\renewcommand{\Aone}{\mathbf A^1}
\providecommand{\Bl}{}\renewcommand{\Bl}{\operatorname{Bl}}
\providecommand{\Br}{}\renewcommand{\Br}{\mathrm{Br}}
\providecommand{\C}{}\renewcommand{\C}{\mathbb C}
\providecommand{\Cech}{}\renewcommand{\Cech}{\check{H}}
\providecommand{\Coext}{}\renewcommand{\Coext}{\operatorname{Coext}}
\providecommand{\Cone}{}\renewcommand{\Cone}{\operatorname{Cone}}
\providecommand{\Cons}{}\renewcommand{\Cons}{\operatorname{Cons}}
\providecommand{\Cref}{}\renewcommand{\Cref}[1]{\text{#1}}
\providecommand{\D}{}\renewcommand{\D}{D}
\providecommand{\DLAFT}{}\renewcommand{\DLAFT}{\operatorname{FT}}
\providecommand{\DLAFr}{}\renewcommand{\DLAFr}{\operatorname{Fr}}
\providecommand{\DLAQ}{}\renewcommand{\DLAQ}{{\overline{\mathbf Q}_\ell}}
\providecommand{\DLRQ}{}\renewcommand{\DLRQ}{{\overline{\mathbf Q}_\ell}}
\providecommand{\Dbc}{}\renewcommand{\Dbc}{D^b_c}
\providecommand{\Dc}{}\renewcommand{\Dc}{D_{\mathrm{cons}}}
\providecommand{\Dctf}{}\renewcommand{\Dctf}{D^b_{ctf}}
\providecommand{\Dic}{}\renewcommand{\Dic}{D_{\IndCons}}
\providecommand{\Dplus}{}\renewcommand{\Dplus}{D^{+}}
\providecommand{\End}{}\renewcommand{\End}{\operatorname{End}}
\providecommand{\Et}{}\renewcommand{\Et}{\mathrm{\acute{e}t}}
\providecommand{\Ex}{}\renewcommand{\Ex}{\mathrm{Ex}}
\providecommand{\Ext}{}\renewcommand{\Ext}{\operatorname{Ext}}
\providecommand{\F}{}\renewcommand{\F}{\mathbb F}
\providecommand{\FT}{}\renewcommand{\FT}{\operatorname{FT}}
\providecommand{\Fbar}{}\renewcommand{\Fbar}{{\overline{\mathbf F}_q}}
\providecommand{\Fell}{}\renewcommand{\Fell}{\mathbf F_\ell}
\providecommand{\Fib}{}\renewcommand{\Fib}{\operatorname{Fib}}
\providecommand{\Fix}{}\renewcommand{\Fix}{\operatorname{Fix}}
\providecommand{\Fp}{}\renewcommand{\Fp}{{\mathbf F_{p}}}
\providecommand{\Fq}{}\renewcommand{\Fq}{{\mathbf F_q}}
\providecommand{\Fqn}{}\renewcommand{\Fqn}{{\mathbf F_{q^{n}}}}
\providecommand{\Frob}{}\renewcommand{\Frob}{\operatorname{Frob}}
\providecommand{\GL}{}\renewcommand{\GL}{\operatorname{GL}}
\providecommand{\Gal}{}\renewcommand{\Gal}{\operatorname{Gal}}
\providecommand{\Hens}{}\renewcommand{\Hens}{\operatorname{Hens}}
\providecommand{\Hilb}{}\renewcommand{\Hilb}{\operatorname{Hilb}}
\providecommand{\Hom}{}\renewcommand{\Hom}{\operatorname{Hom}}
\providecommand{\IndCons}{}\renewcommand{\IndCons}{\operatorname{IndCons}}
\providecommand{\Kar}{}\renewcommand{\Kar}{\operatorname{Kar}}
\providecommand{\LT}{}\renewcommand{\LT}{\operatorname{LT}}
\providecommand{\Lie}{}\renewcommand{\Lie}{\operatorname{Lie}}
\providecommand{\Loc}{}\renewcommand{\Loc}{\operatorname{Loc}}
\providecommand{\Mod}{}\renewcommand{\Mod}{\mathsf{Mod}}
\providecommand{\N}{}\renewcommand{\N}{\mathbb{N}}
\providecommand{\Needspace}{}\renewcommand{\Needspace}{\par}
\providecommand{\O}{}\renewcommand{\O}{{\mathcal{O}}}
\providecommand{\Ob}{}\renewcommand{\Ob}{\mathrm{Ob}}
\providecommand{\Oe}{}\renewcommand{\Oe}{\mathcal O_E}
\providecommand{\PP}{}\renewcommand{\PP}{\mathbf{P}}
\providecommand{\Perf}{}\renewcommand{\Perf}{\operatorname{Perf}}
\providecommand{\Pic}{}\renewcommand{\Pic}{\mathrm{Pic}}
\providecommand{\Pm}{}\renewcommand{\Pm}{\mathbf P^{m}}
\providecommand{\Pone}{}\renewcommand{\Pone}{\mathbf P^{1}}
\providecommand{\Pqn}{}\renewcommand{\Pqn}{\mathbf P^{n}}
\providecommand{\ProFin}{}\renewcommand{\ProFin}{\operatorname{ProFin}}
\providecommand{\Proet}{}\renewcommand{\Proet}{\mathrm{pro\acute{e}t}}
\providecommand{\Q}{}\renewcommand{\Q}{\mathbb Q}
\providecommand{\Qb}{}\renewcommand{\Qb}{\overline{\mathbf Q}_{\ell}}
\providecommand{\Qbar}{}\renewcommand{\Qbar}{\overline{\mathbf Q}}
\providecommand{\QbarX}{}\renewcommand{\QbarX}{\overline{\mathbf Q}_{\ell,X}}
\providecommand{\Qbarl}{}\renewcommand{\Qbarl}{{\overline{\mathbf Q}_{\ell}}}
\providecommand{\Qell}{}\renewcommand{\Qell}{{\mathbf{Q}_{\ell}}}
\providecommand{\Ql}{}\renewcommand{\Ql}{\mathbf Q_{\ell}}
\providecommand{\Qlb}{}\renewcommand{\Qlb}{{\overline{\mathbf Q}_\ell}}
\providecommand{\Qlbar}{}\renewcommand{\Qlbar}{\overline{\mathbf Q}_\ell}
\providecommand{\R}{}\renewcommand{\R}{\mathbb{R}}
\providecommand{\RG}{}\renewcommand{\RG}{\operatorname{\mathrm{R}\Gamma}}
\providecommand{\RGamma}{}\renewcommand{\RGamma}{R\Gamma}
\providecommand{\RGammac}{}\renewcommand{\RGammac}{R\Gamma_c}
\providecommand{\RHom}{}\renewcommand{\RHom}{R\mathcal{H}om}
\providecommand{\RHomX}{}\renewcommand{\RHomX}{R\mathcal H om_X}
\providecommand{\Rat}{}\renewcommand{\Rat}{\operatorname{IsRat}}
\providecommand{\Res}{}\renewcommand{\Res}{\operatorname{Res}}
\providecommand{\Rlim}{}\renewcommand{\Rlim}{R\!\varprojlim}
\providecommand{\SIm}{}\renewcommand{\SIm}{\operatorname{SIm}}
\providecommand{\Sch}{}\renewcommand{\Sch}{\mathsf{Sch}}
\providecommand{\Sh}{}\renewcommand{\Sh}{\operatorname{Sh}}
\providecommand{\Span}{}\renewcommand{\Span}{\operatorname{span}}
\providecommand{\Spec}{}\renewcommand{\Spec}{\operatorname{Spec}}
\providecommand{\Tr}{}\renewcommand{\Tr}{\operatorname{Tr}}
\providecommand{\Trsm}{}\renewcommand{\Trsm}{\operatorname{Tr}^{\mathrm{sm}}}
\providecommand{\Wt}{}\renewcommand{\Wt}{\operatorname{Wt}}
\providecommand{\Z}{}\renewcommand{\Z}{\mathbb{Z}}
\providecommand{\Zar}{}\renewcommand{\Zar}{\mathrm{zar}}
\providecommand{\Zell}{}\renewcommand{\Zell}{{\mathbf{Z}_{\ell}}}
\providecommand{\Zl}{}\renewcommand{\Zl}{{\mathbf Z_\ell}}
\providecommand{\affet}{}\renewcommand{\affet}[1]{{#1}^{\mathrm{aff}}_{\mathrm{ét}}}
\providecommand{\affproet}{}\renewcommand{\affproet}[1]{{#1}^{\mathrm{aff}}_{\mathrm{proét}}}
\providecommand{\boxt}{}\renewcommand{\boxt}{\boxtimes_A^{\mathrm L}}
\providecommand{\cF}{}\renewcommand{\cF}{\mathcal{F}}
\providecommand{\cHom}{}\renewcommand{\cHom}{R\mathcal{H}\!om}
\providecommand{\cI}{}\renewcommand{\cI}{\mathcal{I}}
\providecommand{\cL}{}\renewcommand{\cL}{\mathcal L}
\providecommand{\cM}{}\renewcommand{\cM}{\mathcal M}
\providecommand{\cN}{}\renewcommand{\cN}{\mathcal N}
\providecommand{\cO}{}\renewcommand{\cO}{\mathcal O}
\providecommand{\cT}{}\renewcommand{\cT}{\mathcal T}
\providecommand{\checked}{}\renewcommand{\checked}{\color{red}\aftergroup\restorecolor}
\providecommand{\coeff}{}\renewcommand{\coeff}{\operatorname{coeff}}
\providecommand{\coker}{}\renewcommand{\coker}{\operatorname{coker}}
\providecommand{\colim}{}\renewcommand{\colim}{\operatorname{colim}}
\providecommand{\constructedby}{}\renewcommand{\constructedby}[1]{\par\smallskip\noindent\textcolor{purple!60!black}{\textsf{Constructed by:} }
  \footnotesize\sloppy\nolinkurl{#1}\par}
\providecommand{\deepresult}{}\renewcommand{\deepresult}[1]{\par\smallskip\noindent\textcolor{orange!45!black}{
  \textsf{Blueprint role: named deep theorem node; see the cluster proof architecture.}}
  \par\noindent\textcolor{orange!45!black}{\textsf{Mathematical source:} }#1\par}
\providecommand{\dt}{}\renewcommand{\dt}{\mathbin{\otimes^{\mathbf L}}}
\providecommand{\dw}{}\renewcommand{\dw}{\operatorname{dw}}
\providecommand{\et}{}\renewcommand{\et}[1]{{#1}_{\mathrm{ét}}}
\providecommand{\ev}{}\renewcommand{\ev}{\operatorname{ev}}
\providecommand{\fib}{}\renewcommand{\fib}{\operatorname{fib}}
\providecommand{\formalizationnote}{}\renewcommand{\formalizationnote}[1]{\begin{quote}\small\raggedright\textbf{Formalization note.} #1\end{quote}}
\providecommand{\fpqc}{}\renewcommand{\fpqc}{X_{\mathrm{fpqc}}}
\providecommand{\gammacont}{}\renewcommand{\gammacont}{\Gamma_{\mathrm{cont}}}
\providecommand{\gm}{}\renewcommand{\gm}{\mathbf G_m}
\providecommand{\hocolim}{}\renewcommand{\hocolim}{\operatorname{hocolim}}
\providecommand{\holim}{}\renewcommand{\holim}{\operatorname{holim}}
\providecommand{\hrB}{}\renewcommand{\hrB}{\@ifnextchar]{\hrEnd}{\hrStep}}
\providecommand{\hrEnd}{}\renewcommand{\hrEnd}{\@firstoftwo{}}
\providecommand{\hrStep}{}\renewcommand{\hrStep}{\@firstoftwo{\hrB}}
\providecommand{\id}{}\renewcommand{\id}{\operatorname{id}}
\providecommand{\kunFl}{}\renewcommand{\kunFl}{\mathbf F_\ell}
\providecommand{\kunZn}{}\renewcommand{\kunZn}{\mathbf Z/n\mathbf Z}
\providecommand{\kuncd}{}\renewcommand{\kuncd}{\operatorname{cd}}
\providecommand{\kundim}{}\renewcommand{\kundim}{\dim}
\providecommand{\kunotimes}{}\renewcommand{\kunotimes}[1]{\otimes^{\mathbf L}_{#1}}
\providecommand{\kunsupp}{}\renewcommand{\kunsupp}{\operatorname{Supp}}
\providecommand{\lotimes}{}\renewcommand{\lotimes}{\otimes^{\mathrm L}}
\providecommand{\m}{}\renewcommand{\m}{{\mathfrak{m}}}
\providecommand{\n}{}\renewcommand{\n}{{\mathfrak{n}}}
\providecommand{\nagataGap}{}\renewcommand{\nagataGap}[1]{\par\noindent\textbf{Formalization gap.} #1\par}
\providecommand{\openobligation}{}\renewcommand{\openobligation}[1]{\par\smallskip\noindent\textcolor{red!65!black}{\textsf{Open obligation:} }
  \footnotesize\sloppy\nolinkurl{#1}\par}
\providecommand{\p}{}\renewcommand{\p}{{\mathfrak{p}}}
\providecommand{\pr}{}\renewcommand{\pr}{\operatorname{pr}}
\providecommand{\preshv}{}\renewcommand{\preshv}[1]{\mathrm{PreShv}(#1)}
\providecommand{\pro}{}\renewcommand{\pro}[1]{\mathrm{Pro}(#1)}
\providecommand{\proet}{}\renewcommand{\proet}{X_{\mathrm{proét}}}
\providecommand{\q}{}\renewcommand{\q}{{\mathfrak{q}}}
\providecommand{\rad}{}\renewcommand{\rad}{\operatorname{rad}}
\providecommand{\rank}{}\renewcommand{\rank}{\operatorname{rank}}
\providecommand{\red}{}\renewcommand{\red}{\color{red}}
\providecommand{\restorecolor}{}\renewcommand{\restorecolor}{\color{black}}
\providecommand{\retiredobligation}{}\renewcommand{\retiredobligation}[1]{\par\smallskip\noindent\textcolor{red!75!black}{\textsf{Retired invalid route:} }
  \footnotesize\sloppy\nolinkurl{#1}\par}
\providecommand{\rk}{}\renewcommand{\rk}{\operatorname{rk}}
\providecommand{\sep}{}\renewcommand{\sep}{\textnormal{sep}}
\providecommand{\shv}{}\renewcommand{\shv}[1]{\mathrm{Shv}(#1)}
\providecommand{\spec}{}\renewcommand{\spec}[1]{\mathrm{Spec}(#1)}
\providecommand{\stacks}{}\renewcommand{\stacks}[1]{\href{https://stacks.math.columbia.edu/tag/#1}{\textnormal{Stacks, Tag~#1}}}
\providecommand{\stacksproject}{}\renewcommand{\stacksproject}[1]{{\cite[\href{https://stacks.math.columbia.edu/tag/#1}{Tag #1}]{stacks-project}}}
\providecommand{\statusnote}{}\renewcommand{\statusnote}[1]{\begin{quote}\small\raggedright\textbf{Status.} #1\end{quote}}
\providecommand{\tildeideal}{}\renewcommand{\tildeideal}[2]{\widetilde{(#1, #2)}}
\providecommand{\tildering}{}\renewcommand{\tildering}[3]{{#1}_{\widetilde{(#2, #3)}}}
\providecommand{\tr}{}\renewcommand{\tr}{\operatorname{tr}}
\providecommand{\wt}{}\renewcommand{\wt}{\operatorname{wt}}

\begin{document}
\pagestyle{plain}
\begin{center}
{\LARGE\bfseries Weil Conjectures\par}
\vspace{0.7em}
{\large Appendix F: Algebraic groups in characteristic zero\par}
\vspace{0.5em}
{\small\color{bpmuted}Snapshot 2026-09-26, revision 8e4cb88e2a\par}
\end{center}
\vspace{0.5em}
{\small\tableofcontents}
\bigskip\hrule\bigskip

% object: chap-Algebraic-groups-in-characteristic-zero (blueprint/chap-Algebraic-groups-in-characteristic-zero/_section.md)
\section*{\texorpdfstring{F\quad Algebraic groups in characteristic zero}{F Algebraic groups in characteristic zero}}
\phantomsection\label{bp1}%
\addcontentsline{toc}{section}{\texorpdfstring{F\quad Algebraic groups in characteristic zero}{F Algebraic groups in characteristic zero}}
% text: prose-F-01
\phantomsection\label{bp2}%
Rankin's argument (Chapter \bpout{14}) bounds the eigenvalues of a real sheaf by the largest determinant weight of its constituents, that is, the weight of the determinant divided by the rank. For this it needs the determinant weights of the factors of tensor and exterior powers of a lisse sheaf (Proposition \bpout{14.19}), before anyone knows how these powers decompose. The tool is an element of the arithmetic monodromy group, of positive degree, that acts by a scalar on each simple factor of every tensor construction. This scalar then gives the determinant weight.

In a simple example Frobenius itself does this. Take \(V=\overline{\mathbf Q}_\ell\oplus\overline{\mathbf Q}_\ell(-1)\) on a curve over \(\mathbf F_q\), on which a Frobenius element \(s\) of degree one acts by \(\operatorname{diag}(1,q)\). Here the geometric monodromy is trivial, so the arithmetic monodromy group \(G\) is the torus \(\{\operatorname{diag}(1,t)\}\), and \(\rho(s)\) is central in it. Hence on every tensor construction \(\rho(s)\) acts by a scalar on each simple factor. For instance, it acts by \(1\) and \(q\) on the factors of \(V\), of weights \(0\) and \(2\), and by \(q\) on \(\bigwedge^2V\).

In general \(\rho(s)\) is not central, and we must correct a power of it by an element of the geometric monodromy group \(H\). An element of \(H\) alone would not do, since its determinants on these factors are roots of unity. What we want instead is a central element of the form \(h\rho(s)^N\) with \(h\in H\) and \(N>0\).

Such an element exists by Lemma \bpout{14.15}. For a semisimple Weil representation \(V\), the identity component \(G^\circ\) is generated by \(Z(G^\circ)^\circ\) and the identity component \(H^\circ\). So a power of Frobenius lying in \(G^\circ\) is an element of \(Z(G^\circ)^\circ\) times an element of \(H^\circ\). Averaging over the components of \(G\) then makes it central in \(G\). This is the algebraic form of a familiar fact about compact Lie groups: a connected compact group is the product of the identity component of its center and its derived group.

For algebraic groups over \(\overline{\mathbf Q}_\ell\) we argue with the Lie algebra \(\mathfrak g\) of \(G^\circ\). Since \(\mathfrak g\) acts faithfully and semisimply on \(V\), it is reductive, that is,

\[
\mathfrak g=Z(\mathfrak g)+[\mathfrak g,\mathfrak g]
\]

(Proposition \hyperref[bp28]{F.8}). The derived algebra lies in the Lie algebra of \(H\), because commutators of \(G\) lie in \(H\). It remains to use that a closed connected subgroup of \(G^\circ\) with the full Lie algebra is all of \(G^\circ\).

The tools behind this argument need characteristic zero. Cartier's theorem that algebraic groups are smooth (Proposition \hyperref[bp5]{F.1}) lets us recognize a closed subgroup by its Lie algebra; this fails for \(\mu_p\subset\mathbf G_m\) in characteristic \(p\). Similarly, the decomposition of \(\mathfrak g\) rests on Cartan's criterion, which fails there too. Since the monodromy groups are defined over \(\overline{\mathbf Q}_\ell\), this does not restrict us. For the theorems of Lie, Engel and Cartan we refer to Humphreys [\hyperref[bp30]{Hum72}].

Throughout, \(\Omega\) is an algebraically closed field of characteristic zero, and all Lie algebras and representations are finite-dimensional. By an algebraic group we mean a group scheme of finite type over \(\Omega\). In particular, stabilizers and centralizers are taken as subgroup schemes, and by Cartier's theorem below they are smooth. The Lie algebra \(\operatorname{Lie}(G)\) of an algebraic group \(G\) is its tangent space at the identity; for a closed subgroup of \(\GL(V)\) we describe it by dual numbers in Section \hyperref[bp10]{F.2}.

% object: sec-Smoothness-and-components (blueprint/chap-Algebraic-groups-in-characteristic-zero/sec-Smoothness-and-components/_section.md), content hidden
\subsection*{\texorpdfstring{F.1\quad Smoothness and components}{F.1 Smoothness and components}}
\phantomsection\label{bp3}%
\addcontentsline{toc}{subsection}{\texorpdfstring{F.1\quad Smoothness and components}{F.1 Smoothness and components}}
% text: prose-F-02
\phantomsection\label{bp4}%
To pass from Lie algebras back to groups we use the following principle: a closed subgroup of a connected group \(G\) with the same Lie algebra is all of \(G\) (Lemma \hyperref[bp8]{F.2}). For reduced groups this is a dimension count: the dimension of a smooth group equals that of its tangent space at \(1\), and a closed subset of the irreducible variety \(G\) of the same dimension is all of \(G\). The difficulty is that stabilizers and centralizers are defined by equations and need not be reduced. For example, in characteristic \(p\), the subgroup \(\mu_p\subset\mathbf G_m\) defined by \(x^p=1\) is a single nonreduced point with tangent space \(\operatorname{Lie}(\mathbf G_m)\), and the principle fails. However, Cartier's theorem (Proposition \hyperref[bp5]{F.1}) shows that this cannot happen in characteristic zero: every algebraic group is smooth. In the proof we apply the comultiplication to a nilpotent function at the identity and divide by an integer, and this division is where we use characteristic zero. The facts about connected components and images of homomorphisms in Lemma \hyperref[bp8]{F.2} are standard consequences of smoothness and of Chevalley's theorem.

% object: proposition-cartier-s-theorem (blueprint/chap-Algebraic-groups-in-characteristic-zero/sec-Smoothness-and-components/proposition-cartier-s-theorem/_section.md)
\phantomsection\label{bp5}%
\begin{bpstatementit}{Proposition~F.1}{ (Cartier's theorem)}
Every algebraic group over \(\Omega\) is smooth and reduced. In particular

\[
\dim\operatorname{Lie}(G)=\dim G.
\]
\end{bpstatementit}

% proof: proof/proposition-cartier-s-theorem
\begin{proof}[{Proof}]\phantomsection\label{bp6}%
We first show, using translations, that the reduced subscheme \(G_{\mathrm{red}}\) is a smooth group scheme. It then suffices to show that \(G\) has the same tangent space at \(1\) as \(G_{\mathrm{red}}\), that is, that nilpotent functions at \(1\) lie in \(\mathfrak m^2\), and this is where we divide by an integer.

\emph{Step 1: \(G_{\mathrm{red}}\) is a smooth subgroup scheme of \(G\).} Multiplication and inversion restrict to \(G_{\mathrm{red}}\). This is because products of reduced schemes of finite type over the perfect field \(\Omega\) are reduced [Stacks, \href{https://stacks.math.columbia.edu/tag/035Z}{035Z}], and a morphism from a reduced scheme to \(G\) factors through \(G_{\mathrm{red}}\). The reduced scheme \(G_{\mathrm{red}}\) contains a dense smooth open subscheme [Stacks, \href{https://stacks.math.columbia.edu/tag/056V}{056V}]. It is smooth everywhere, since a translation moves a closed point of this open to any other closed point of \(G_{\mathrm{red}}\).

\emph{Step 2: nilpotent functions at \(1\) lie in \(\mathfrak m^2\).} Let \(R=\mathcal O_{G,1}\), with maximal ideal \(\mathfrak m\) and nilradical \(I\), and let \(\widehat R\) be its completion. At the point \((1,1)\) the multiplication induces a continuous comultiplication

\[
\Delta:\widehat R\longrightarrow
\widehat R\,\widehat\otimes_\Omega\widehat R.
\]

The counit is evaluation at \(1\), and composing \(\Delta\) with evaluation at \(1\) in either factor gives the identity (the counit identities), because \(1\) is a two-sided unit for the group law. We claim that \(I\subseteq\mathfrak m^2\). As completion is faithfully flat, it suffices to prove that every nilpotent \(b\in\widehat R\) lies in \(\widehat{\mathfrak m}^{\,2}\). Suppose that \(b\notin\widehat{\mathfrak m}^{\,2}\), and let \(n\ge2\) be the least integer with \(b^n=0\). By the counit identities,

\[
\Delta(b)=b\otimes1+1\otimes b+w,
\qquad
w\in\widehat{\mathfrak m}\,\widehat\otimes_\Omega
\widehat{\mathfrak m}.
\]

Now reduce the second factor modulo \(\widehat{\mathfrak m}^{\,2}\). Every term in the binomial expansion of \(\Delta(b)^n\) with at least two factors from \(1\otimes b+w\) vanishes, because \(1\otimes b+w\) has its second factor in \(\widehat{\mathfrak m}\), and as \(b^n=0\) we are left with

\[
n b^{n-1}\otimes b+n(b^{n-1}\otimes1)w=0.
\]

If we now also reduce the first factor modulo \(b^{n-1}\widehat{\mathfrak m}\), the second term vanishes, and we get

\[
n\,\overline{b^{n-1}}\otimes\overline b=0
\quad\text{in}\quad
(\widehat R/b^{n-1}\widehat{\mathfrak m})
\otimes_\Omega(\widehat R/\widehat{\mathfrak m}^{\,2}).
\]

This identity is impossible, which proves the claim. Indeed, the class \(\overline{b^{n-1}}\) is nonzero by Nakayama's lemma, because the module \((b^{n-1})\) is nonzero by the choice of \(n\). The class \(\overline b\) is nonzero by assumption. A tensor product of nonzero vectors over a field is nonzero, and \(n\ne0\) in characteristic zero.

\emph{Step 3: \(G\) is smooth.} The schemes \(G\) and \(G_{\mathrm{red}}\) have the same cotangent space at \(1\). Indeed, by Step 2 the kernel of \(\mathfrak m/\mathfrak m^2\to(\mathfrak m/I)/(\mathfrak m/I)^2\) is \((I+\mathfrak m^2)/\mathfrak m^2=0\). They also have the same local dimension at \(1\), because a nilpotent ideal does not change the underlying space. As the local ring \(R/I\) of \(G_{\mathrm{red}}\) at \(1\) is regular by Step 1, we get

\[
\dim_\Omega\mathfrak m/\mathfrak m^2=\dim R,
\]

so \(R\) is regular, and in particular a domain [Stacks, \href{https://stacks.math.columbia.edu/tag/00NP}{00NP}]. By translation the local rings of \(G\) at all closed points are regular. Since \(\Omega\) is perfect, \(G\) is smooth at every closed point. The nonsmooth locus is closed and contains no closed point, and hence it is empty. Being smooth, \(G\) is reduced, and its tangent space at \(1\) has dimension \(\dim G\). Finally, the open subgroup \(G^\circ\) has the same tangent space at \(1\).
\end{proof}

% text: prose-F-03
\phantomsection\label{bp7}%
Since algebraic groups over \(\Omega\) are smooth, their topology is as simple as that of a Lie group: irreducible and connected components coincide, and the dimension of a closed subgroup can be computed from its Lie algebra.

% object: lemma-connected-components-images-homomorphisms (blueprint/chap-Algebraic-groups-in-characteristic-zero/sec-Smoothness-and-components/lemma-connected-components-images-homomorphisms/_section.md)
\phantomsection\label{bp8}%
\begin{bpstatementit}{Lemma~F.2}{ (Connected components and images of homomorphisms)}
Let \(G\) be an algebraic group.

\begin{enumerate}
\item \(G\) has finitely many connected components. Its identity component \(G^\circ\) is an open, closed, irreducible normal subgroup, the other components are its cosets, and \(\operatorname{Lie}(G^\circ)=\operatorname{Lie}(G)\).
\item The image of a homomorphism \(G\to G'\) of algebraic groups is a closed subgroup of \(G'\).
\item If \(G\) is connected and \(K\subseteq G\) is a closed subgroup scheme with \(\operatorname{Lie}(K)=\operatorname{Lie}(G)\), then \(K=G\).
\end{enumerate}
\end{bpstatementit}

% proof: proof/lemma-connected-components-images-homomorphisms
\begin{proof}[{Proof}]\phantomsection\label{bp9}%
For (1) we use smoothness to see that irreducible and connected components coincide; for (2), Chevalley's theorem gives an open subset of the closure inside the image, and translations spread it over the closure; (3) is a dimension count.

\emph{Components.} The irreducible components of \(G\) are finite in number, as \(G\) is Noetherian, and they are exactly the connected components. Indeed, \(G\) is smooth by Proposition \hyperref[bp5]{F.1}, and a regular local ring is a domain. So distinct irreducible components of \(G\) do not meet, and they are open and closed. Let \(G^\circ\) be the component containing \(1\). It is a subgroup. To see this, note that \(G^\circ\times G^\circ\) is irreducible, \(\Omega\) being algebraically closed. Hence the image of the morphism \(G^\circ\times G^\circ\to G\), \((x,y)\mapsto xy^{-1}\), is irreducible. It contains \(1\), so it lies in \(G^\circ\). For the same reason the image of \(G^\circ\) under conjugation by a point of \(G\) is irreducible and contains \(1\), so \(G^\circ\) is normal. Translating \(G^\circ\) by points of \(G\), we see that the other components are its cosets. Since \(G^\circ\) is open in \(G\) and contains \(1\), both have the same local ring at \(1\), so \(\operatorname{Lie}(G^\circ)=\operatorname{Lie}(G)\).

For example, the homomorphism \(\mathbf G_m\to\mathbf G_m\), \(x\mapsto x^2\), has image all of \(\mathbf G_m\) and kernel \(\mu_2\), a closed subgroup with two components.

\emph{Images.} Let \(f:G\to G'\) be a homomorphism. We show that its image equals the closure \(J\) of the image, which is a subgroup by continuity of multiplication and inversion. First, the subgroup \(f(G(\Omega))\) of \(J(\Omega)\) contains the \(\Omega\)-points of a nonempty open \(U\subseteq J\). Indeed, \(f(G)\) is constructible by Chevalley's theorem [Stacks, \href{https://stacks.math.columbia.edu/tag/054K}{054K}], a dense constructible subset contains a nonempty open subset of its closure, and \(\Omega\) is algebraically closed. Translating \(U\) by the inverse of one of its points in the image, we obtain an open neighbourhood of \(1\) contained in the image. Translating further by points of the image, we see that the image is open in \(J\). It is then also closed, since its complement is a union of cosets of this open subgroup. As it is dense in \(J\), it is equal to \(J\).

\emph{Lie algebras.} Let \(K\) be as in (3). Cartier's theorem gives \(\dim K=\dim\operatorname{Lie}(K)=\dim G\), so \(K\) and \(G\) have the same underlying closed subset, as a proper closed subset of the irreducible variety \(G\) has smaller dimension. They are equal as schemes, since both are reduced by Proposition \hyperref[bp5]{F.1}.
\end{proof}


% object: sec-Lie-algebras (blueprint/chap-Algebraic-groups-in-characteristic-zero/sec-Lie-algebras/_section.md), content hidden
\subsection*{\texorpdfstring{F.2\quad Lie algebras}{F.2 Lie algebras}}
\phantomsection\label{bp10}%
\addcontentsline{toc}{subsection}{\texorpdfstring{F.2\quad Lie algebras}{F.2 Lie algebras}}
% text: prose-F-04
\phantomsection\label{bp11}%
Let \(G\) be a connected algebraic group with Lie algebra \(\mathfrak g\). We need two facts (Lemma \hyperref[bp12]{F.3}). First, in a representation, \(G\) and \(\mathfrak g\) have the same stable subspaces. Second, the Lie algebra of the identity component of the center of \(G\) is the center of \(\mathfrak g\). In the proof of Lemma \bpout{14.15} we compare the monodromy groups through their Lie algebras, and both facts are used there. By the first, a representation that is semisimple for \(G^\circ\) is semisimple for its Lie algebra. This is what lets Proposition \hyperref[bp28]{F.8} apply to the Weil representation. With the second we translate the decomposition of the Lie algebra into the identity \(G^\circ=Z(G^\circ)^\circ H^\circ\).

We prove both facts by the principle of the previous section, applied to a stabilizer and to the kernel of the adjoint representation. Their Lie algebras are computed with dual numbers. The first fact needs \(G\) to be connected. For example, the finite group \(\{\pm1\}\subset\GL_1\) has zero Lie algebra, so every subspace of a representation is stable under its Lie algebra. But if \(-1\) acts on \(M=\Omega^2\) by \(\operatorname{diag}(1,-1)\), then it does not stabilize the line spanned by \((1,1)\).

For a closed subgroup \(G\subseteq\GL(V)\), the tangent space \(\operatorname{Lie}(G)\subseteq\operatorname{End}(V)\) is described by dual-number points at \(1\):

\[
X\in\operatorname{Lie}(G)
\quad\Longleftrightarrow\quad
1+\varepsilon X\in G(\Omega[\varepsilon]/\varepsilon^2).
\]

The bracket is \([X,Y]=XY-YX\). To see that \(\operatorname{Lie}(G)\) is closed under the bracket, note that over \(\Omega[\varepsilon,\eta]/(\varepsilon^2,\eta^2)\) the commutator of \(1+\varepsilon X\) and \(1+\eta Y\) is \(1+\varepsilon\eta[X,Y]\). Its entries lie in the subring \(\Omega[\varepsilon\eta]\), which is a ring of dual numbers. The equations of \(G\) hold for this matrix in the larger ring, hence in this subring. Therefore \([X,Y]\in\operatorname{Lie}(G)\). The differential of a representation \(r:G\to\GL(M)\) is characterized by

\[
r(1+\varepsilon X)=1+\varepsilon\,dr(X).
\]

Applying \(r\) to the commutator of \(1+\varepsilon X\) and \(1+\eta Y\), we see that \(dr[X,Y]=[dr(X),dr(Y)]\), so \(dr\) is a Lie algebra homomorphism.

% object: lemma-stable-subspaces-center (blueprint/chap-Algebraic-groups-in-characteristic-zero/sec-Lie-algebras/lemma-stable-subspaces-center/_section.md)
\phantomsection\label{bp12}%
\begin{bpstatementit}{Lemma~F.3}{ (Stable subspaces and the center)}
Let \(G\subseteq\GL(V)\) be a connected closed subgroup and \(r:G\to\GL(M)\) an algebraic representation. A subspace of \(M\) is \(G\)-stable if and only if it is \(\operatorname{Lie}(G)\)-stable. Moreover

\[
\operatorname{Lie}\bigl(Z(G)^\circ\bigr)
=Z\bigl(\operatorname{Lie}(G)\bigr).
\]
\end{bpstatementit}

% proof: proof/lemma-stable-subspaces-center
\begin{proof}[{Proof}]\phantomsection\label{bp13}%
In both assertions we compare a closed subgroup with \(G\), namely the stabilizer of a subspace in the first and the kernel of the adjoint representation in the second. We compute the Lie algebra of this subgroup with dual numbers and conclude with Lemma lemma-connected-components-images-homomorphisms.

Let \(U\subseteq M\) be a subspace. In a basis of \(M\) extending a basis of \(U\), the scheme-theoretic stabilizer \(S\subseteq G\) of \(U\) is the closed subgroup defined by the vanishing of the lower-left block of \(r(g)\). Since the lower-left block of \(r(1+\varepsilon X)=1+\varepsilon\,dr(X)\) is \(\varepsilon\) times that of \(dr(X)\), we find

\[
\operatorname{Lie}(S)
=\{X\in\operatorname{Lie}(G):dr(X)U\subseteq U\}.
\]

If \(U\) is \(\operatorname{Lie}(G)\)-stable, then \(\operatorname{Lie}(S)=\operatorname{Lie}(G)\), and Lemma lemma-connected-components-images-homomorphisms shows that \(S=G\). Conversely, if \(U\) is \(G\)-stable, then \(S=G\), and the displayed formula shows that \(U\) is \(\operatorname{Lie}(G)\)-stable.

Put \(\mathfrak g=\operatorname{Lie}(G)\). We claim that the kernel of the adjoint representation \(G\to\GL(\mathfrak g)\) is the center. Indeed, if \(g\) acts trivially on \(\mathfrak g\), consider the centralizer \(C_G(g)\). It is a closed subgroup scheme whose Lie algebra is the fixed space of \(\operatorname{Ad}(g)\), that is, all of \(\mathfrak g\). So \(C_G(g)=G\) by Lemma lemma-connected-components-images-homomorphisms. Conversely, a central element acts trivially. Since the kernel and the center are reduced by Cartier's theorem, their equality on \(\Omega\)-points is an equality of schemes.

The Lie algebra of the center is now the kernel of a differential. The formula

\[
\operatorname{Ad}(1+\varepsilon X)(Y)=Y+\varepsilon[X,Y]
\]

shows that the differential of \(\operatorname{Ad}\) is \(\operatorname{ad}\). The center is a closed subgroup of \(\GL(V)\), so the dual-number description applies to it. By this description, the Lie algebra of the kernel of a representation is the kernel of its differential. Hence \(\operatorname{Lie}(Z(G))=\ker(\operatorname{ad})=Z(\mathfrak g)\). By Proposition \hyperref[bp5]{F.1}, \(Z(G)^\circ\) has the same Lie algebra as \(Z(G)\).
\end{proof}


% object: sec-Cartans-criterion (blueprint/chap-Algebraic-groups-in-characteristic-zero/sec-Cartans-criterion/_section.md), content hidden
\subsection*{\texorpdfstring{F.3\quad Cartan's criterion}{F.3 Cartan's criterion}}
\phantomsection\label{bp14}%
\addcontentsline{toc}{subsection}{\texorpdfstring{F.3\quad Cartan's criterion}{F.3 Cartan's criterion}}
% text: prose-F-05
\phantomsection\label{bp15}%
The decomposition \(\mathfrak l=Z(\mathfrak l)+[\mathfrak l,\mathfrak l]\) fails for Lie algebras in general. Consider the algebra \(\mathfrak b\) of upper triangular \(2\times2\) matrices. Its center consists of the scalars, and \([\mathfrak b,\mathfrak b]\) consists of the strictly upper triangular matrices. So their sum has dimension two and misses \(\operatorname{diag}(1,0)\). Here \(\mathfrak b\) is solvable and is its own radical, much larger than its center. Its standard representation on \(\Omega^2\) is faithful but not semisimple, since the line \(\Omega e_1\) has no stable complement. In general the obstruction is the radical, that is, the largest solvable ideal. In Proposition \hyperref[bp28]{F.8} we show that a faithful semisimple representation forces the radical to be central. This gives the decomposition.

The derived series of a Lie algebra is \(D^0\mathfrak l=\mathfrak l\), \(D^{n+1}\mathfrak l=[D^n\mathfrak l,D^n\mathfrak l]\), and \(\mathfrak l\) is solvable if this series reaches zero. The radical \(\operatorname{rad}\mathfrak l\) is the sum of all solvable ideals of \(\mathfrak l\). To see that it is itself solvable, note that an extension of solvable Lie algebras is solvable. Indeed, if the quotient of \(\mathfrak l\) by a solvable ideal is solvable, the derived series of \(\mathfrak l\) lies in that ideal after finitely many steps, and then reaches zero. Hence the sum of two solvable ideals is solvable, and \(\operatorname{rad}\mathfrak l\) is the largest solvable ideal of \(\mathfrak l\) [\hyperref[bp30]{Hum72}, §3.1]. From the same book we use three classical theorems, and we only sketch their proofs.

% object: lemma-lie-s-theorem (blueprint/chap-Algebraic-groups-in-characteristic-zero/sec-Cartans-criterion/lemma-lie-s-theorem/_section.md)
\phantomsection\label{bp16}%
\begin{bpstatementit}{Lemma~F.4}{ (Lie's theorem)}
A solvable Lie algebra acting on a nonzero vector space \(M\) has a common eigenvector. Equivalently, there are \(v\ne0\) and a linear functional \(\lambda\) such that \(av=\lambda(a)v\) for every \(a\).
\end{bpstatementit}

% text: prose-F-06
\phantomsection\label{bp17}%
See [\hyperref[bp30]{Hum72}, §4.1]. By induction one finds a common eigenvector for an ideal of codimension one, and then shows by a trace computation, which we redo in the proof of Proposition \hyperref[bp28]{F.8}, that the corresponding eigenspace is stable under the remaining direction.

% object: lemma-engel-s-theorem (blueprint/chap-Algebraic-groups-in-characteristic-zero/sec-Cartans-criterion/lemma-engel-s-theorem/_section.md)
\phantomsection\label{bp18}%
\begin{bpstatementit}{Lemma~F.5}{ (Engel's theorem)}
Let \(\mathfrak n\subseteq\operatorname{End}(M)\) be a Lie algebra consisting of nilpotent endomorphisms. If \(M\ne0\), then the common kernel of \(\mathfrak n\) is nonzero. There is a basis of \(M\) in which all elements of \(\mathfrak n\) are strictly upper triangular. In particular, \(\mathfrak n\) is nilpotent and solvable.
\end{bpstatementit}

% text: prose-F-07
\phantomsection\label{bp19}%
This is Engel's theorem in its matrix form [\hyperref[bp30]{Hum72}, §3.3]. The first assertion is proved by induction on \(\dim\mathfrak n\), using that a maximal proper subalgebra is an ideal of codimension one. The triangular basis then follows by induction on \(\dim M\).

% object: lemma-cartan-s-criterion (blueprint/chap-Algebraic-groups-in-characteristic-zero/sec-Cartans-criterion/lemma-cartan-s-criterion/_section.md)
\phantomsection\label{bp20}%
\begin{bpstatementit}{Lemma~F.6}{ (Cartan's criterion)}
Let \(M\) be a finite-dimensional \(\Omega\)-vector space, and let \(\mathfrak a\subseteq\operatorname{End}(M)\) be a Lie subalgebra satisfying

\[
\operatorname{tr}(XY)=0
\qquad(X\in[\mathfrak a,\mathfrak a],\ Y\in\mathfrak a).
\]

Then \(\mathfrak a\) is solvable.
\end{bpstatementit}

% text: prose-F-08
\phantomsection\label{bp21}%
See [\hyperref[bp30]{Hum72}, §4.3]. It suffices to show that every element \(X\) of \([\mathfrak a,\mathfrak a]\) is nilpotent. Then \([\mathfrak a,\mathfrak a]\) is solvable by Lemma \hyperref[bp18]{F.5}, and hence so is \(\mathfrak a\). To see that \(X\) is nilpotent, one pairs \(X\) with operators that are polynomials in its semisimple part. Let \(\lambda_i\) be the eigenvalues of the semisimple part, and let \(f\) be a \(\mathbf Q\)-linear functional on their \(\mathbf Q\)-span with values in \(\mathbf Q\). The hypothesis gives \(\sum_i\lambda_i f(\lambda_i)=0\). Applying \(f\) to this identity gives \(\sum_i f(\lambda_i)^2=0\). This is a vanishing sum of squares of rational numbers, so \(f\) vanishes on all \(\lambda_i\), and \(f=0\). As \(f\) was arbitrary, all \(\lambda_i\) are zero. This step is where we use characteristic zero.


% object: sec-Reductive-Lie-algebras (blueprint/chap-Algebraic-groups-in-characteristic-zero/sec-Reductive-Lie-algebras/_section.md), content hidden
\subsection*{\texorpdfstring{F.4\quad Reductive Lie algebras}{F.4 Reductive Lie algebras}}
\phantomsection\label{bp22}%
\addcontentsline{toc}{subsection}{\texorpdfstring{F.4\quad Reductive Lie algebras}{F.4 Reductive Lie algebras}}
% text: prose-F-09
\phantomsection\label{bp23}%
For \(\mathfrak{gl}_n\) the decomposition into scalars and trace-zero matrices, \(\mathfrak{gl}_n=\Omega\cdot1\oplus\mathfrak{sl}_n\), has the form \(\mathfrak l=Z(\mathfrak l)+[\mathfrak l,\mathfrak l]\). We prove this decomposition for every Lie algebra with a faithful semisimple representation (Proposition \hyperref[bp28]{F.8}), as Lemma \bpout{14.15} requires. The proof has two steps. First, the radical acts by scalars on each simple summand. To see this, we take a common eigenvector given by Lie's theorem. A trace computation, in which we divide by the dimension of a subspace, shows that its eigenspace is stable under all of \(\mathfrak l\). Since the representation is faithful, the radical is then central. Second, the quotient by the center has zero radical. For such a Lie algebra the Killing form is nondegenerate by Cartan's criterion. This forces the algebra to equal its derived algebra (Lemma \hyperref[bp24]{F.7}). The upper triangular algebra \(\mathfrak b\) of the previous section shows that the hypothesis of semisimplicity cannot be dropped.

% object: lemma-killing (blueprint/chap-Algebraic-groups-in-characteristic-zero/sec-Reductive-Lie-algebras/lemma-killing/_section.md)
\phantomsection\label{bp24}%
\begin{bpstatementit}{Lemma~F.7}{ (Killing form)}
If \(\operatorname{rad}\mathfrak l=0\), the Killing form \(B(x,y)=\operatorname{tr}(\operatorname{ad}x\operatorname{ad}y)\) is nondegenerate and \([\mathfrak l,\mathfrak l]=\mathfrak l\).
\end{bpstatementit}

% text: prose-F-10
\phantomsection\label{bp25}%
This is the semisimplicity criterion of [\hyperref[bp30]{Hum72}, §§5.1–5.2]; we include the short deduction from Cartan's criterion.

% proof: proof/lemma-killing
\begin{proof}[{Proof of Lemma F.7}]\phantomsection\label{bp26}%
We first use Cartan's criterion to show that the kernel of \(B\) is a solvable ideal, hence zero; nondegeneracy then identifies the orthogonal of the derived algebra with the center. By invariance of \(B\), that is, \(B([x,y],z)=B(x,[y,z])\), the kernel \(\mathfrak k\) of \(B\) is an ideal. For \(x\in[\mathfrak k,\mathfrak k]\) and \(y\in\mathfrak k\), the operators \(\operatorname{ad}x\) and \(\operatorname{ad}y\) map \(\mathfrak l\) into \(\mathfrak k\), so the trace of \(\operatorname{ad}x\operatorname{ad}y\) may be computed on \(\mathfrak k\). Since \(x\in\mathfrak k\), we obtain

\[
\operatorname{tr}_{\mathfrak k}
(\operatorname{ad}(x)|_{\mathfrak k}
\operatorname{ad}(y)|_{\mathfrak k})=B(x,y)=0.
\]

By Cartan's criterion (Lemma \hyperref[bp20]{F.6}), the image of \(\mathfrak k\) in \(\operatorname{End}(\mathfrak k)\) is solvable. The kernel of \(\mathfrak k\to\operatorname{End}(\mathfrak k)\) is the center of \(\mathfrak k\), which is abelian. So \(\mathfrak k\) is a solvable ideal of \(\mathfrak l\). Since \(\operatorname{rad}\mathfrak l=0\), we get \(\mathfrak k=0\), that is, \(B\) is nondegenerate. Invariance and nondegeneracy then give

\[
[\mathfrak l,\mathfrak l]^\perp=Z(\mathfrak l).
\]

The center is a solvable ideal, hence zero, and a subspace with zero orthogonal complement is all of \(\mathfrak l\).
\end{proof}

% text: prose-F-11
\phantomsection\label{bp27}%
For \(\mathfrak{sl}_2\), with basis \(e,h,f\), the brackets \([h,e]=2e\), \([h,f]=-2f\) and \([e,f]=h\) already show that \([\mathfrak{sl}_2,\mathfrak{sl}_2]=\mathfrak{sl}_2\). The lemma gives the second half of the next proposition, once we know that the quotient by the center has zero radical.

% object: proposition-lie-algebras-faithful-semisimple-representation (blueprint/chap-Algebraic-groups-in-characteristic-zero/sec-Reductive-Lie-algebras/proposition-lie-algebras-faithful-semisimple-representation/_section.md)
\phantomsection\label{bp28}%
\begin{bpstatementit}{Proposition~F.8}{ (Lie algebras with a faithful semisimple representation)}
If \(\mathfrak l\) has a faithful semisimple representation, then

\[
\operatorname{rad}\mathfrak l=Z(\mathfrak l),
\qquad
\mathfrak l=Z(\mathfrak l)+[\mathfrak l,\mathfrak l].
\]
\end{bpstatementit}

% proof: proof/proposition-lie-algebras-faithful-semisimple-representation
\begin{proof}[{Proof}]\phantomsection\label{bp29}%
We follow the two halves described at the start of this section: first the radical is central, and then the quotient by the center equals its derived algebra.

Put \(\mathfrak r=\operatorname{rad}\mathfrak l\), and write the faithful module as a direct sum of simple modules. Let \(M\) be one of them. By Lemma \hyperref[bp16]{F.4}, there is a nonzero simultaneous \(\lambda\)-eigenspace \(M_\lambda\) for \(\mathfrak r\). We claim that \(M_\lambda\) is \(\mathfrak l\)-stable. To prove this, fix \(x\in\mathfrak l\) and \(0\ne v\in M_\lambda\), and consider \(U=\operatorname{span}\{v,xv,x^2v,\ldots\}\). For every \(a\in\mathfrak r\) we have \([a,x]\in\mathfrak r\), and the identity

\[
ax^j=x^ja+\sum_{i=0}^{j-1}x^i[a,x]x^{j-1-i}
\]

shows, by induction on \(j\), that \(a\) preserves \(U\) and acts on it by a triangular matrix with constant diagonal \(\lambda(a)\) in the basis \(v,xv,x^2v,\ldots\) of \(U\). Indeed, the identity gives \(ax^jv=\lambda(a)x^jv\) plus terms \(x^i[a,x]x^{j-1-i}v\), and by induction applied to \([a,x]\in\mathfrak r\) these lie in the span of \(v,\ldots,x^{j-1}v\). With this, the trace of \([x,a]\) on \(U\) can be computed in two ways. Since \(x\) and \(a\) preserve \(U\), the restriction of \([x,a]\) to \(U\) is a commutator, whose trace is zero. On the other hand, \([x,a]\in\mathfrak r\) acts on \(U\) with constant diagonal \(\lambda([x,a])\). Comparing,

\[
0=\operatorname{tr}_U[x,a]=\dim(U)\lambda([x,a]),
\]

and \(\lambda([x,a])=0\) because \(\Omega\) has characteristic zero. It follows that \(a(xv)=\lambda(a)xv+\lambda([a,x])v=\lambda(a)xv\), which proves the claim.

Since \(M\) is simple, the claim gives \(M_\lambda=M\). Thus \(\mathfrak r\) acts on \(M\) by scalars, and in particular \([\mathfrak l,\mathfrak r]\) acts as zero on \(M\). This holds for every simple summand, and the module is faithful, so \([\mathfrak l,\mathfrak r]=0\), that is, \(\mathfrak r\subseteq Z(\mathfrak l)\). The reverse inclusion holds because the center is an abelian ideal.

For the second equality, we show that the quotient \(\overline{\mathfrak l}=\mathfrak l/Z(\mathfrak l)\) has zero radical and apply Lemma \hyperref[bp24]{F.7}. The inverse image of a solvable ideal of the quotient is a central extension of a solvable Lie algebra, hence solvable. So it is contained in \(\mathfrak r=Z(\mathfrak l)\). The lemma now gives \([\overline{\mathfrak l},\overline{\mathfrak l}]
 =\overline{\mathfrak l}\). In other words, every element of \(\mathfrak l\) is congruent modulo \(Z(\mathfrak l)\) to a sum of brackets, and \(\mathfrak l=Z(\mathfrak l)+[\mathfrak l,\mathfrak l]\).
\end{proof}



\clearpage\phantomsection\hypertarget{bp-outside}{}%
\section*{Outside this export}
\addcontentsline{toc}{section}{Outside this export}
References marked \ensuremath{{}^\ast} are to parts of the blueprint this document leaves out:
\begin{description}
\item[Chapter 14] Purity of real sheaves, in Ⅲ Weights and the Riemann hypothesis (\bpslug{chap-{}\allowbreak{}Purity-{}\allowbreak{}of-{}\allowbreak{}real-{}\allowbreak{}sheaves})
\item[Lemma 14.15] The identity component, in Ⅲ Weights and the Riemann hypothesis (\bpslug{lemma-{}\allowbreak{}identity-{}\allowbreak{}component})
\item[Proposition 14.19] Determinant weights of tensor operations, in Ⅲ Weights and the Riemann hypothesis (\bpslug{proposition-{}\allowbreak{}determinant-{}\allowbreak{}weights-{}\allowbreak{}tensor-{}\allowbreak{}operations})
\end{description}
\clearpage\phantomsection%
\section*{References}
\addcontentsline{toc}{section}{References}
\bpbibitem{BBD82}{bp31}{A. Beilinson, J. Bernstein, and P. Deligne. \emph{Faisceaux pervers}, volume 100 of \emph{Astérisque}. 1982.}

\bpbibitem{BS15}{bp32}{B. Bhatt and P. Scholze. The pro-étale topology for schemes. \emph{Astérisque}, 369:99–201, 2015. \href{https://people.mpim-bonn.mpg.de/scholze/proetale.pdf}{Author's manuscript}.}

\bpbibitem{Con07}{bp33}{B. Conrad. Deligne's notes on Nagata compactifications. \emph{Journal of the Ramanujan Mathematical Society}, 22:205–257, 2007. \href{https://math.stanford.edu/\~conrad/papers/nagatafinal.pdf}{Author's manuscript}.}

\bpbibitem{Del74}{bp34}{P. Deligne. La conjecture de Weil. I. \emph{Publ. Math. IHÉS}, 43:273–307, 1974.}

\bpbibitem{Del80}{bp35}{P. Deligne. La conjecture de Weil. II. \emph{Publ. Math. IHÉS}, 52:137–252, 1980.}

\bpbibitem{Den22}{bp36}{C. Deninger. There is no “Weil-”cohomology theory with real coefficients for arithmetic curves. \emph{Ann. Sc. Norm. Super. Pisa Cl. Sci.}, 2022. To appear. arXiv:\href{https://arxiv.org/abs/2204.02714}{2204.02714}.}

\bpbibitem{Deu41}{bp37}{M. Deuring. Die Typen der Multiplikatorenringe elliptischer Funktionenkörper. \emph{Abh. Math. Sem. Univ. Hamburg}, 14:197–272, 1941.}

\bpbibitem{Eke90}{bp38}{T. Ekedahl. On the adic formalism. In \emph{The Grothendieck Festschrift, Vol. II}, volume 87 of \emph{Progress in Mathematics}, pages 197–218. Birkhäuser, 1990.}

\bpbibitem{Has36}{bp39}{H. Hasse. Zur Theorie der abstrakten elliptischen Funktionenkörper. I, II, III. \emph{J. Reine Angew. Math.}, 175:55–62, 69–88, 193–208, 1936.}

\bpbibitem{Hum72}{bp30}{J. E. Humphreys. \emph{Introduction to Lie Algebras and Representation Theory}, volume 9 of \emph{Graduate Texts in Mathematics}. Springer, 1972.}

\bpbibitem{Ill94}{bp40}{L. Illusie. Autour du théorème de monodromie locale. \emph{Astérisque}, 223:9–57, 1994. §4. \href{https://www.numdam.org/item/AST_1994__223__9_0/}{Original article}.}

\bpbibitem{Kat88}{bp41}{N. M. Katz. Travaux de Laumon. \emph{Astérisque}, 161–162:105–132, 1988. Séminaire Bourbaki, exp. 691.}

\bpbibitem{KW01}{bp42}{R. Kiehl and R. Weissauer. \emph{Weil conjectures, perverse sheaves and \(\ell\)-adic Fourier transform}, volume 42 of \emph{Ergebnisse der Mathematik}. Springer, 2001.}

\bpbibitem{Lau87}{bp43}{G. Laumon. Transformation de Fourier, constantes d'équations fonctionnelles et conjecture de Weil. \emph{Publ. Math. IHÉS}, 65:131–210, 1987.}

\bpbibitem{Mil86a}{bp44}{J. S. Milne. \href{https://www.jmilne.org/math/xnotes/AVs.pdf}{Abelian Varieties}. In \emph{Arithmetic Geometry}. 1986. Corrected online version.}

\bpbibitem{Mil86b}{bp45}{J. S. Milne. \href{https://www.jmilne.org/math/xnotes/JVs.pdf}{Jacobian Varieties}. In \emph{Arithmetic Geometry}. 1986. Corrected online version.}

\bpbibitem{Ran39}{bp46}{R. A. Rankin. Contributions to the theory of Ramanujan's function \(\tau(n)\) and similar arithmetical functions. II. \emph{Proc. Cambridge Philos. Soc.}, 35:357–372, 1939.}

\bpbibitem{SGA4}{bp47}{M. Artin, A. Grothendieck, and J.-L. Verdier, editors. \emph{Théorie des topos et cohomologie étale des schémas}, volumes 269, 270, 305 of \emph{Lecture Notes in Mathematics}. Springer, 1972–1973. SGA 4.}

\bpbibitem{SGA4½}{bp48}{P. Deligne. \emph{Cohomologie étale}, volume 569 of \emph{Lecture Notes in Mathematics}. Springer, 1977. SGA 4½; \emph{Théorèmes de finitude}, §3, and \emph{Rapport sur la formule des traces}, §3, and \emph{Sommes trigonométriques}, §1.}

\bpbibitem{SGA5}{bp49}{A. Grothendieck, editor. \emph{Cohomologie \(\ell\)-adique et fonctions \(L\)}, volume 589 of \emph{Lecture Notes in Mathematics}. Springer, 1977. SGA 5.}

\bpbibitem{SGA7II}{bp50}{P. Deligne and N. Katz, editors. \emph{Groupes de monodromie en géométrie algébrique}, volume 340 of \emph{Lecture Notes in Mathematics}. Springer, 1973. SGA 7 II; Exposé XIII, (2.1.7).}

\bpbibitem{Stacks}{bp51}{The Stacks Project Authors. \href{https://stacks.math.columbia.edu}{The Stacks Project}. Individual results are cited by their stable tags in the text.}

\bpbibitem{Var05}{bp52}{Y. Varshavsky. Lefschetz–Verdier trace formula and a generalization of a theorem of Fujiwara. 2005. Version 2. arXiv:\href{https://arxiv.org/abs/math/0505564}{math/0505564}. \href{https://arxiv.org/abs/math/0505564v2}{Author's preprint}.}

\bpbibitem{Wei48a}{bp53}{A. Weil. \emph{Sur les courbes algébriques et les variétés qui s'en déduisent}. Hermann, Paris, 1948.}

\bpbibitem{Wei48b}{bp54}{A. Weil. \emph{Variétés abéliennes et courbes algébriques}. Hermann, Paris, 1948.}

\bpbibitem{Wei49}{bp55}{A. Weil. Numbers of solutions of equations in finite fields. \emph{Bull. Amer. Math. Soc.}, 55:497–508, 1949.}

\end{document}
